The empirical results highlight a significant divergence in model performance between the Global Financial Crisis (GFC) and the COVID-19 pandemic. During the GFC, yield curve factors—specifically Slope ($S_t$) and Curvature ($C_t$)—remained highly informative because the recession was preceded by a classic restrictive monetary environment where the term premium eroded1111. This is consistent with the foundational work of Estrella and Mishkin (1998), which establishes the yield spread as a primary barometer for financially driven downturns.However, the COVID-19 episode presented a unique challenge for machine learning models. The Curvature factor (PC3), which represents the "hump" or convexity of the curve, showed episodic spikes during this period of intense policy intervention2222. As the Federal Reserve implemented large-scale asset purchases, the "belly" of the curve moved independently of the short and long ends, reflecting market pricing of recovery risk against a backdrop of aggressive stimulus3. This supports the narrative in Rudebusch, Sack, and Swanson (2007) that unconventional monetary policy can compress term premia and distort the pure expectations-driven signals of the yield curve4.The rolling PCA diagnostics provide the technical evidence for this "regime change." Significant drift was observed in the factor loadings during both the GFC and COVID-19 eras, suggesting that the underlying relationship between different maturities is not stationary5. Because the COVID-19 recession was an exogenous shock rather than an endogenous economic overheating, the yield curve factors were "structurally distorted"6. Consequently, macroeconomic indicators and labor market variables dominated forecasting performance during this period, as they captured the immediate reality of the lockdown shock that the forward-looking yield curve failed to anticipate twelve months in advance777.This underscores the "Maschine Learning Basics" mentioned in your goal: while flexible models like XGBoost or Random Forests can capture non-linearities, they remain sensitive to the training regime8. When a "Black Swan" event occurs that is absent from the training data, even a high-performing model will struggle unless it incorporates broader macroeconomic features beyond the yield curve999.